%%%%%%%%%%  scan idx 58  |  folio 47  %%%%%%%%%%

\underline{Lema 1.18. Supongamos que la condición (II$'$) se cumple
para $\underline A$. Entonces:}

(1) $[\underline A(D_{K'}(K))]$ está provisto de una única estructura
de espacio vectorial tal que, funtorialmente en $V$,
\[
 [\underline A(D_V(K))]
 \simeq
 V\otimes_{K'}[\underline A(D_{K'}(K))].
\]

(2) Existe una aplicación lineal
\[
 \gamma:\operatorname{Hom}(\Omega,K')
 \longrightarrow [\underline A(D_{K'}(K))]
\]
tal que, para todo $D\in\operatorname{Hom}(\Omega,K')$, el morfismo
\[
 [\underline A((1,D))]:
 [\underline A(D_{K'}(K))]
 \longrightarrow [\underline A(D_{K'}(K))]
\]
es $x\longmapsto x+\gamma(D)$.

(3) Para todo
$(u,D):D_V(K)\longrightarrow D_W(K)$, el morfismo inducido
\[
 (u,D):[\underline A(D_V(K))]
 \longrightarrow[\underline A(D_W(K))],
\]
es decir,
\[
 (u,D):
 V\otimes[\underline A(D_{K'}(K))]
 \longrightarrow
 W\otimes[\underline A(D_{K'}(K))],
\]
viene dado por
\[
 x\longmapsto u(x)+\gamma(D).
\]

Introduzcamos la notación
\[
\begin{aligned}
 F(V)&=[\underline A(D_V(K))],\\
 A&=F(K'),\\
 F(u,D):F(V)&\longrightarrow F(W),\\[-.25ex]
 &\quad\text{para el morfismo inducido por}\\[-.25ex]
 &\quad (u,D):D_V(K)\longrightarrow D_W(K),\\
 F(D):F(V)&\longrightarrow F(V),
   &F(D)&=F(1,D).
\end{aligned}
\]
Por (II$'$), y puesto que $K'$ es un objeto espacio vectorial sobre
$K'$ en la categoría de los espacios vectoriales sobre $K'$,
$F(K')$ es un objeto espacio vectorial sobre $K'$ en la categoría de
los conjuntos y, por tanto, un espacio vectorial sobre $K'$. Esto
demuestra (1). En adelante identificaremos $F(V)$ con $V\otimes A$.

La identidad
\[
 (u,D)=(1,D)\circ(u,0)
\]
da
\[
 F(u,D)=F(D)\circ(u\otimes1_A).
\]
La funtorialidad de $F$ equivale a las identidades
\[
\begin{aligned}
 F(0)&=\operatorname{Id},\\
 F(D)\circ F(D')&=F(D+D'),\\
 (u\otimes1)\circ F(D)
   &=F(uD)\circ(u\otimes1_A).
\end{aligned}
\]

Para todo $D:\Omega\longrightarrow V$ y todo $n\in\mathbb N$, sean
\[
 i:V\lhook\joinrel\longrightarrow V\oplus K'^{\,n},
 \qquad
 p:V\oplus K'^{\,n}\longrightarrow \sisi{K^{n}}{K'^{n}}
\]
la inclusión y la proyección canónicas. Entonces $piD=0$, y el
diagrama siguiente es conmutativo:

%%%%%%%%%%  scan idx 59  |  folio 48  %%%%%%%%%%

\[
\begin{tikzcd}[column sep=huge,row sep=large]
 V\otimes A
   \arrow[r]
   \arrow[d,"{F(D)}"]
 & (V\oplus K'^{\,n})\otimes A
   \arrow[r]
   \arrow[d,"{\sisi{F(i,D)}{F(i\,D)}}"]
 & K'^{\,n}\otimes A
   \arrow[d,equal]\\
 V\otimes A
   \arrow[r]
 & (V\oplus K'^{\,n})\otimes A
   \arrow[r]
 & K'^{\,n}\otimes A.
\end{tikzcd}
\]

Sea $x:K'^{\,n}\longrightarrow V$ un morfismo, pongamos
$x_i=x(e_i)$ y sea $(a_i)$ una familia de elementos de $A$. Se tiene
\[
 (1_V+x)\circ i\circ D=D.
\]
Por consiguiente, si
\[
 F(i,D)\Bigl(\sum_i e_i a_i\Bigr)
 =\sum_i e_i a_i+a,
 \qquad a\in V\otimes A,
\]
entonces
\[
 F(D)\Bigl(\sum_i x_i a_i\Bigr)
 =\sum_i x_i a_i+a,
\]
y $a$ es independiente de los $x_i$. Por tanto,
\[
\begin{aligned}
 F(D)(y)&=y+G(D),\\
 G(D+D')&=G(D)+G(D'),\\
 uG(D)&=G(uD)
   \qquad\text{para todo }u:V\longrightarrow W.
\end{aligned}
\]
Todo $G$ de esta forma viene dado, para algún
$\gamma\in\Omega\otimes A$, por
\[
 G(D)=(D\otimes1_A)(\gamma).
\]
Esto demuestra (2) y (3).
